% !TEX root = ../main_proposal.tex
% =====================================================================
%  BAB 1 — PENDAHULUAN
% =====================================================================

\chapter{\MakeUppercase{Pendahuluan}}

\section{Latar Belakang}
\label{sec:1.1}

Emosi merupakan komponen mendasar dalam komunikasi antarmanusia, karena
emosi memungkinkan seseorang menyampaikan maksud dan keadaan batin
secara lebih utuh daripada makna literal kata-kata. Dalam interaksi
manusia dan mesin, pengenalan emosi dari sinyal ucapan memegang peranan
penting bagi pengembangan komputasi afektif dan sistem interaktif yang
mampu merespons pengguna secara lebih personal (Dar and Delhibabu,
2024). Karakteristik inilah yang mendasari berkembangnya \emph{speech
emotion recognition} (SER), yaitu bidang kajian yang berupaya mengenali
kondisi emosi seseorang secara otomatis dari sinyal ucapannya, dengan
cakupan penerapan yang terus meluas. Meskipun signifikansinya diakui
luas, perkembangan teknologi bahasa berbasis ucapan masih timpang
antarbahasa, karena sebagian besar teknologi kecerdasan buatan terpusat
pada sebagian kecil bahasa dengan sumber daya digital terbesar (Bella et
al., 2023). Bahasa Indonesia, meskipun dituturkan ratusan juta jiwa,
termasuk kelompok bahasa yang sumber dayanya masih tertinggal, sehingga
data ucapan beranotasi emosi dalam jumlah memadai masih sulit diperoleh.
Kesenjangan inilah yang menjadikan pengembangan SER berbahasa Indonesia
sebagai permasalahan penelitian yang penting sekaligus menantang.

Menjawab kesenjangan tersebut, Santoso, Budianto and Dutono (2026)
mengembangkan E-SERAVD, dataset SER berbahasa Indonesia yang bersumber
dari dialog rekaman film Indonesia, sekaligus menetapkan \emph{baseline}
performa 86\% akurasi dan 86\% F1-macro menggunakan model CNN
konvensional dengan fitur \emph{hand-crafted} MFCC. \emph{Baseline} tersebut,
meskipun berguna, dilaporkan tanpa menyebutkan kontrol terhadap sumber
film pada pembagian data latih dan uji. Padahal, E-SERAVD terbukti
memiliki sebaran film yang tidak merata antarkelas emosi. Terdapat 106 dari 300
klip kelas sedih berasal dari satu film yang sama, dan 147 dari 300 klip
kelas bahagia berasal dari dua film yang sama. Apabila pembagian data latih dan uji dilakukan secara
acak tanpa mengontrol sumber film, klip dari film yang sama berisiko
muncul pada kedua sisi sekaligus, sehingga model berpotensi mempelajari
ciri produksi atau karakteristik rekaman film tertentu sebagai proxy
kelas emosi, bukan pola akustik emosi secara umum. 

Risiko tersebut
dikenal sebagai kebocoran data (\emph{data leakage}) (Kapoor and
Narayanan, 2023), dan berpotensi membuat performa yang dilaporkan
terlihat lebih baik daripada performa sesungguhnya pada data yang
benar-benar baru. Pada domain SER secara spesifik, kegagalan mengontrol
tumpang tindih identitas penutur antara data latih dan uji telah
terbukti menghasilkan estimasi performa yang lebih optimistis
dibandingkan skema evaluasi yang independen (Antoniou et al., 2023). Pada domain SER
secara umum, kegagalan mengontrol kebocoran semacam ini dilaporkan
memengaruhi sebagian besar penelitian yang diterbitkan. Lebih dari
80\% studi SER tidak dapat direproduksi secara memadai, dan skema
partisi data yang tidak terkontrol dilaporkan menghasilkan peningkatan
performa semu hingga sekitar 4\% (Wu et al., 2024).

Di sisi lain, representasi \emph{self-supervised} kini banyak dianggap
sebagai pendekatan performa terbaik (\emph{state-of-the-art}) untuk SER,
termasuk untuk bahasa Indonesia. Atmaja and Sasou (2022) menunjukkan
bahwa representasi \emph{self-supervised} secara konsisten mengungguli
fitur \emph{hand-crafted} pada lima dataset SER berbahasa Inggris dan Jepang, dengan
signifikansi statistik yang kuat, meskipun
XLSR-53 secara spesifik sempat berperforma lebih rendah dibandingkan
model monolingual berbahasa Inggris pada studi yang sama. Pada konteks
SER berbahasa Indonesia secara spesifik, Nugroho et al. (2025)
menunjukkan bahwa model dasar XLSR-Wav2Vec2 yang telah di-\emph{fine-tune}
untuk ucapan Indonesia dapat mencapai akurasi SER hingga 91,67\%,
mengungguli metode klasik berbasis fitur \emph{hand-crafted}. Bukti-bukti ini
menempatkan XLSR-53 sebagai backbone \emph{self-supervised} yang paling
relevan untuk diadaptasi pada tugas SER berbahasa Indonesia saat ini. Akan tetapi, belum
ada studi yang menguji ketahanan algoritma \emph{baseline} (CNN+MFCC)
maupun algoritma SOTA (XLSR-53 yang diadaptasi dengan LoRA) terhadap
risiko kebocoran data berbasis film yang terbukti ada pada E-SERAVD.

Berdasarkan uraian diatas,
diperlukan kajian analitik mengenai pengaruh skema pembagian data
berbasis film terhadap akurasi pengenalan emosi ucapan berbahasa
Indonesia, serta apakah pengaruh tersebut konsisten pada algoritma
berbasis fitur \emph{hand-crafted} maupun algoritma berbasis representasi
\emph{self-supervised},
yang akan dianalisis dalam skripsi ini.

\section{Rumusan Masalah}
\label{sec:1.2}

Berdasarkan latar belakang yang telah dipaparkan, maka permasalahan
dalam penelitian ini dapat dirumuskan sebagai berikut.

\begin{enumerate}
  \item Bagaimana pengaruh skema pembagian data berbasis film
        (\emph{random} dibandingkan \emph{movie-independent}) pada dataset E-SERAVD terhadap
        performa klasifikasi pengenalan emosi ucapan berbahasa
        Indonesia?
  \item Apakah pengaruh tersebut konsisten pada algoritma berbasis
        fitur \emph{hand-crafted} (MFCC+CNN) maupun algoritma berbasis representasi
        \emph{self-supervised} (XLSR-53)?
\end{enumerate}

\section{Tujuan}
\label{sec:1.3}

Tujuan dari penelitian ini menjawab rumusan masalah yang telah
dipaparkan, maka penelitian ini memiliki tujuan sebagai berikut.

\begin{enumerate}
  \item Mengetahui pengaruh skema pembagian data berbasis film
        (\emph{random} dibandingkan \emph{movie-independent}) terhadap
        performa klasifikasi pengenalan emosi ucapan berbahasa
        Indonesia.
  \item Menganalisis konsistensi pengaruh skema pembagian data berbasis
        film (\emph{random split} dibandingkan \emph{movie-independent
        split}) terhadap akurasi pengenalan emosi ucapan berbahasa
        Indonesia pada algoritma berbasis fitur \emph{hand-crafted} (MFCC+CNN) dan
        algoritma berbasis representasi \emph{self-supervised} (XLSR-53).
\end{enumerate}

\section{Manfaat}
\label{sec:1.4}

Penelitian ini diharapkan dapat memberikan manfaat sebagai berikut.

\begin{enumerate}
  \item Memberikan bukti empiris mengenai validitas \emph{baseline}
        performa yang telah dilaporkan pada dataset E-SERAVD terhadap
        risiko kebocoran data, sehingga dapat menjadi rekomendasi
        metodologis bagi peneliti lain yang menggunakan dataset ini
        pada penelitian selanjutnya.
  \item Memberikan bukti empiris mengenai ketahanan algoritma berbasis
        representasi \emph{self-supervised} dibandingkan algoritma berbasis fitur
        \emph{hand-crafted} terhadap skema pembagian data yang lebih ketat, sehingga
        dapat membantu peneliti maupun praktisi dalam memilih algoritma
        yang lebih andal untuk SER berbahasa Indonesia.
  \item Berkontribusi pada pengembangan riset komputasi cerdas berbahasa
        Indonesia, khususnya pada bidang pemrosesan sinyal ucapan yang
        masih relatif terbatas jumlah kajiannya.
\end{enumerate}

\section{Batasan Masalah}
\label{sec:1.5}

Batasan masalah yang membatasi penulis dalam melakukan penelitian ini
adalah sebagai berikut.

\begin{enumerate}
  \item Dataset dibatasi pada E-SERAVD v1.1, mencakup 1.200 klip ucapan
        emosional berbahasa Indonesia dari 6 film sumber, dengan 4
        kelas emosi (marah, sedih, netral, bahagia).
  \item Algoritma \emph{baseline} dibatasi pada CNN dengan fitur \emph{hand-crafted}
        MFCC, mengikuti pendekatan Santoso, Budianto and Dutono (2026).
  \item Algoritma pembanding dibatasi pada backbone XLSR-53 yang diadaptasi
        menggunakan LoRA (\emph{Low-Rank Adaptation}), bukan
        \emph{full fine-tuning} pada seluruh parameter backbone.
  \item Skema pembagian data dibatasi pada dua skema: \emph{random} dan
        \emph{movie-independent}, keduanya menggunakan $k=6$
        \emph{fold}.
  \item Evaluasi performa dibatasi pada metrik kualitas klasifikasi yaitu
        akurasi, presisi, recall, dan F1-score.
  \item Signifikansi statistik diuji menggunakan uji Shapiro-Wilk untuk
        normalitas, dilanjutkan uji \emph{paired t-test} atau
        \emph{Wilcoxon signed-rank test} sesuai hasil uji normalitas
        tersebut.
\end{enumerate}

\section{Sistematika Pembahasan}
\label{sec:1.6}

\textbf{Bab 1 Pendahuluan}

Membahas mengenai latar belakang penelitian
ini diangkat, rumusan masalah, tujuan dan manfaat penelitian, batasan
masalah, serta sistematika pembahasan penulisan penelitian yang
berjudul Pengaruh Skema Pembagian Data Berbasis Film pada Dataset E-SERAVD terhadap Performa Klasifikasi
Pengenalan Emosi Ucapan Berbahasa Indonesia.

\textbf{Bab 2 Landasan Kepustakaan}

Membahas tentang kajian pustaka dan
dasar teori. Kajian pustaka berisi referensi mengenai
penelitian-penelitian sebelumnya yang terkait dengan SER dan kebocoran
data pada skema evaluasi. Dasar teori membahas konsep-konsep yang
mendasari metode yang digunakan.

\textbf{Bab 3 Metodologi Penelitian}

Membahas tipe penelitian serta
tahapan penelitian yang dilaksanakan mulai dari studi literatur,
pengumpulan data, perancangan algoritma, evaluasi k-fold, hingga
analisis hasil dan penarikan kesimpulan.

\textbf{Bab 4 Perancangan}

Membahas perancangan sistem serta algoritme
yang digunakan dalam eksperimen penelitian ini, meliputi rancangan
sistem klasifikasi untuk algoritma \emph{baseline} dan algoritma pembanding,
rancangan keempat kondisi eksperimen, serta rancangan skema
validasi silang k-fold.

\textbf{Bab 5 Implementasi}

Membahas implementasi kode-kode esensial
dari perancangan yang telah diuraikan pada Bab sebelumnya, meliputi implementasi
ekstraksi fitur, pelatihan model, dan pengujian signifikansi statistik.

\textbf{Bab 6 Hasil dan Pembahasan}

Menyajikan hasil pelaksanaan keempat
kondisi eksperimen, meliputi observasi akurasi, presisi, recall, dan
F1-score per fold, beserta hasil pengujian normalitas dan signifikansi
statistiknya, serta pembahasan yang menerjemahkan makna hasil tersebut
untuk menjawab rumusan masalah penelitian.

\textbf{Bab 7 Penutup}

Memuat kesimpulan yang menjawab rumusan masalah
berdasarkan hasil dan pembahasan penelitian, serta saran untuk
pengembangan penelitian selanjutnya.
